\chapter{Epilogue}
\label{ch:epilogue}

For the main chapters of this book, I pretended that we were all born with eight fingers, and so had natively adopted a base-8 world---except for a few footnotes where I was comparing a concept to its base-10 cousin. For most people, arithmetic skill and early math intuition are built with numbers written in that native world. Once that intuition is in place, we move on to algebra, geometry, trigonometry, calculus, and so on, where the actual number representation no longer matters that much. So I genuinely had a question: would a native base-8 world be better for learning math?

For a few things, it is actually pretty nice. The multiplication table is smaller. Counting the non-trivial pairs in the table, you only need $21$ of them, against $36$ in a base-10 world. That is nearly a half reduction. It has to help early on: I have watched a lot of kids struggle with the ``large'' $9 \times 9$ table. The fact that $8$ is a cube of $2$ is also nice. Powers of two pick up a trailing zero every three steps. A leading $1$ in Benford's law has chance exactly $1/3$, and one cycle splits into three equal doublings. And a fraction has a finite decimal-point form exactly when the denominator is a power of $2$---which itself is easy to spot: the number starts with $1$, $2$, or $4$, and every digit after that is $0$. It must be a blast to learn binary in a computer science class: every octal digit is just three bits.

Simplicity has an opposite side, though. Although deciding whether a fraction has a finite decimal form is much cleaner in a base-8 world, it is also a bit boring. In our base-10 world, $1/2$, $1/4$, $1/5$, $1/8$, $1/10$, $1/16$, $1/20$, and $1/25$ all have a nice finite decimal-point form. The factors $2$ and $5$ share the work of making tens: a number ends with a zero only when both appear, and a decimal terminates only when the denominator's primes are among those two. Even $100$ itself shows the difference. In base-10 its factors are $1$, $2$, $4$, $5$, $10$, $20$, $25$, $50$, and $100$---powers of $2$, powers of $5$, and products of both, from $2^{2} \times 5^{2}$. In base-8 the factors of $100$ are just $1$, $2$, $4$, $10$, $20$, $40$, and $100$, every one a power of $2$, because $100 = 2^{6}$---quite boring. The variety in the base-10 world makes the math feel mysterious and interesting. Discovering the underlying rule behind a seeming mystery is satisfying, and that feeling can spark more joy in learning math.

Besides the boringness, some arithmetic does seem harder to grasp in a base-8 world. The noticeable one is divisibility by $5$, a very small prime, and one you need for factorization, common denominators, and the like. There does not seem to be a shortcut: kids in a base-8 world will probably just have to memorize the multiples of $5$. Take recognizing two-digit primes as an example. In a base-10 world, the tests for $2$, $3$, and $5$ are so easy that you can spot those multiples at once. If you know the single-digit multiplication table, multiples of $7$ up through $70$ are familiar, $77$ is obvious, and $84$ is even so you never need the $7$-test there. The only non-trivial two-digit composite you really have to memorize is $91 = 7 \times 13$; everything else falls back on ordinary intuition. In a base-8 world, divisibility by $3$ is a bit more complex in general, but for a two-digit number it is simple: check whether the difference of the two digits is a multiple of $3$. So $25$ is, and $12$ is not. Divisibility by $7$ is just the digit sum, so that is pretty trivial too. Running a similar argument, most two-digit composites are not hard to recognize. Assuming you memorized $31$ as a multiple of $5$ in the multiplication table, the only non-trivial two-digit composite number you have to memorize in a base-8 world is $67$, which is a non-obvious multiple of $5$. So is it harder to learn primes in a base-8 world? I do not know. Maybe slightly harder, but not a lot harder.

So I am not trying to crown a winner for early math education. Base-8 makes some patterns inevitable and some tables smaller. Base-10 keeps a mix of $2$ and $5$ that makes a few puzzles more interesting. There is no clear winner for me here, and I just hope you found this book entertaining. It is a pretty good coffee-table book, isn't it? Not serious stuff, but interesting enough for some nerd talk.
